\section{What a direction call is worth}
\label{sec:exp:econ}

Balanced accuracy says how often a model is right; it says nothing about what being right is
worth once the call is traded, because a prediction market charges a spread and, on some
markets, a fee, and because a resting order is filled only when somebody trades against it.
This section asks what a direction call of the models of this chapter is worth under those
costs.  A simple model of a traded call comes first; a light simulation on the locked test
days then measures its terms and the profit or loss of three ways of executing the calls,
as a taker and as a maker.

\subsection{A simple model of a traded call}
\label{sec:econ:model}

A \emph{call} is a sample $(m, t)$ at which a model predicts \textsc{up} or \textsc{down};
the \textsc{const} predictions are not traded.  Trading one share of the Yes token on a call
means buying it on \textsc{up} and selling it on \textsc{down} (which is the same as buying
one share of No at one minus the price), and closing the position at $t + H$.  Three things
can happen by then: the mid has moved in the predicted direction, with probability $h$ (the
\emph{precision} of the calls) and by $\mu^{+}$ on average; it has not moved, with
probability $z$; or it has moved the other way, with probability $w = 1 - h - z$ and by
$\mu^{-}$ on average.  Before any cost the call is worth, per share,
\begin{equation}
  g \;=\; h\,\mu^{+} - w\,\mu^{-},
  \label{eq:econ:edge}
\end{equation}
the \emph{gross edge}: the mean signed move of the mid over the horizon, taken in the
direction of the call.  Balanced accuracy does not fix $g$: it averages the recall of the
three classes, while $g$ depends on the precision of the calls that are actually made and on
the size of the moves, which the label does not see.  Two models with the same balanced
accuracy can have different gross edges, and the same model has a very different $g$ on a
market that moves by one tick and on one that moves by fifty.

The costs depend on how the call is executed.  A \emph{taker} crosses the book twice: it
buys at the ask and later sells at the bid (or the reverse), so it pays the half spread at
entry and the half spread at exit, $s$ in all when the spread does not change, plus a fee on
each leg.  The fee of Polymarket on a leg of $C$ shares at price $p$ is
\begin{equation}
  f \;=\; C\,\rho\,p\,(1 - p),
  \label{eq:econ:fee}
\end{equation}
where $\rho$ is the rate of the market; only takers pay it, makers never do.  In March 2026
the rate was about $6.25\%$ on the short-lived crypto markets and $1\%$ on a few others, and
zero on most markets of the dataset: $72\%$ of the taker legs of the period carry a fee,
almost all of them on crypto markets outside the order-book archive, and among the markets
of the benchmark only the crypto ones pay it on most of their fills.  At $p = 0.5$ the fee is
at most $\rho / 4$ per share, that is $15.6$ ticks at the highest rate and $2.5$ ticks at
$1\%$; at $p = 0.05$ or $0.95$ it is five times smaller.  The net value of a taker's call is
therefore
\begin{equation}
  \mathbb{E}[\text{P\&L}_{\text{taker}}] \;=\; g - s - f_{\text{in}} - f_{\text{out}},
  \label{eq:econ:taker}
\end{equation}
and since $g \le (1 - z)\,\mu^{+}$ even for a model that is never wrong, a taker can only
profit where the spread is smaller than the expected move: $(1 - z)\,\mu^{+} > s + 2f$ is a
necessary condition that no precision can repair.  With symmetric moves
$\mu^{+} = \mu^{-} = \mu$ the break-even precision of a taker is
\begin{equation}
  h^{*} \;=\; \frac{1 - z}{2} + \frac{s + 2f}{2\mu},
  \label{eq:econ:breakeven}
\end{equation}
which exceeds one as soon as the costs exceed the move.

A \emph{maker} rests an order at the best quote, a bid for \textsc{up} and an ask for
\textsc{down}: it pays no spread, earns half of it relative to the mid, and pays the fee only
if it closes as a taker.  Its cost is of another kind.  The order is filled only when a
counterparty crosses to its price, with some probability $q$ over the horizon, and the fills
are not a random sample of the calls: a bid is hit when the price is coming down, so the
gross edge conditional on a fill, $g_{\text{filled}}$, is smaller than $g$ and can be
negative (\emph{adverse selection}).  The value of a maker's call is
$q\,(g_{\text{filled}} + s/2 - f_{\text{out}})$ when the position is closed by crossing, and
$q\,(g_{\text{filled}} + s)$ if the closing order also rests and is filled, an upper bound
that assumes a second fill.  Which side is better is an empirical question, and the
simulation answers it.

\subsection{A light simulation on the test days}
\label{sec:econ:simulation}

The simulation trades every call of a model on the test days of 24--26 March, one share per
call, under the following rules.  A sample is traded only if its book is well formed (a
positive spread, both quotes inside the unit interval; $96\%$ of the samples).  The calls
are those of the calibrated model, with the offsets chosen on the validation day exactly as
in \cref{sec:horizons:test}: the stateless control at every horizon, the stateful model with
the engineered features and its residual on the trees where their test predictions exist,
and an \emph{oracle} that calls the realised label, the upper bound of any direction model
under the same costs.  Three executions are compared.  \emph{Taker--taker} crosses the
spread at $t$ and again at $t + H$, at the quotes of the last book seen before $t + H$
around the true mid of the label.  \emph{Maker--taker} rests at the best quote at $t$; the
order counts as filled if a later fill of the same market within $(t, t + H]$ printed at or
through its price, and the position is closed by crossing at $t + H$; calls that are not
filled are not traded.  \emph{Maker--maker} rests on both legs and assumes the closing fill:
the upper bound of a pure maker.  Fees are those charged in March 2026 on each market, read
from its own taker fills, with a no-fee variant for comparison.  Our own orders do not move
the book, the queue position is ignored (a resting order is filled as soon as the price
reaches it), there is no latency and no capital constraint: the taker numbers are realistic
for one share, the maker numbers are optimistic, and the gap between the two is where a real
strategy would sit.

\Cref{tab:econ:terms} gives the terms of \cref{eq:econ:edge,eq:econ:taker} for the calls of
the control at each horizon, and \cref{tab:econ:pnl} the profit or loss per share of the
three executions; \cref{app:results:econ} has the same quantities for the stateful models,
the oracle and the subsets.

\begin{table}[t]
\centering
\caption{The terms of a traded call for the stateless control on the test days: calls
(millions), precision $h$, share of calls with no move $z$ and with the opposite move $w$,
mean absolute move when right ($\mu^{+}$) and when wrong ($\mu^{-}$), gross edge $g$, median
spread at the calls, mean spread paid by a taker over the two legs, and mean fee of the two
legs as charged in March 2026.  Ticks of $0.001$.}
\label{tab:econ:terms}
\small
\begin{fittable}\begin{tabular}{lrccccccccc}
\toprule
Horizon & Calls & $h$ & $z$ & $w$ & $\mu^{+}$ & $\mu^{-}$ & $g$ & Spread (median) & Spread paid & Fee \\
\midrule
1 min & 5.80 & 0.270 & 0.578 & 0.152 & 17.6 & 19.1 & $+1.9$ & 30 & 100.7 & 0.81 \\
15 min & 8.63 & 0.398 & 0.372 & 0.230 & 32.7 & 37.2 & $+4.5$ & 27 & 95.7 & 0.64 \\
1 h & 9.97 & 0.464 & 0.266 & 0.270 & 43.2 & 43.8 & $+8.2$ & 23 & 89.3 & 0.55 \\
6 h & 9.73 & 0.539 & 0.135 & 0.326 & 33.3 & 31.2 & $+7.8$ & 20 & 81.2 & 0.38 \\
12 h & 9.01 & 0.561 & 0.105 & 0.334 & 35.2 & 31.5 & $+9.2$ & 20 & 77.1 & 0.37 \\
\bottomrule
\end{tabular}\end{fittable}
\end{table}

\begin{table}[t]
\centering
\caption{Profit or loss per share, in ticks of $0.001$, of the calls of the stateless
control on the test days under the three executions, fees as charged in March 2026 unless
stated.  Taker--taker on all calls, without fees, and on the calls with a spread of at most
two ticks; maker--taker on all calls and on the most confident decile, with the share of
calls that are filled; maker--maker on all calls (upper bound); the oracle that calls the
realised label, as a taker and as a maker--taker.}
\label{tab:econ:pnl}
\small
\begin{fittable}\begin{tabular}{lccccccccccc}
\toprule
& \multicolumn{3}{c}{Taker--taker} & \multicolumn{2}{c}{Maker--taker, all} & \multicolumn{2}{c}{Maker--taker, top decile} & Maker--maker & \multicolumn{3}{c}{Oracle} \\
\cmidrule(lr){2-4}\cmidrule(lr){5-6}\cmidrule(lr){7-8}\cmidrule(lr){9-9}\cmidrule(lr){10-12}
Horizon & all & no fee & spread $\le 2$ & filled & P\&L & filled & P\&L & all & taker & filled & maker--taker \\
\midrule
1 min & $-99.6$ & $-98.8$ & $-0.9$ & 7\% & $-4.2$ & 7\% & $+0.6$ & $+24.5$ & $-90.6$ & 10\% & $+27.4$ \\
15 min & $-91.8$ & $-91.2$ & $+0.9$ & 20\% & $+0.4$ & 19\% & $+51.1$ & $+33.0$ & $-63.9$ & 24\% & $+71.5$ \\
1 h & $-81.6$ & $-81.0$ & $+0.8$ & 29\% & $+11.1$ & 45\% & $+79.4$ & $+44.0$ & $-47.5$ & 31\% & $+82.1$ \\
6 h & $-73.8$ & $-73.4$ & $-0.3$ & 46\% & $+10.4$ & 49\% & $+46.2$ & $+40.0$ & $-48.0$ & 41\% & $+40.0$ \\
12 h & $-68.2$ & $-67.8$ & $+0.8$ & 52\% & $+13.9$ & 55\% & $+59.2$ & $+42.6$ & $-42.5$ & 48\% & $+39.1$ \\
\bottomrule
\end{tabular}\end{fittable}
\end{table}

\paragraph{The terms.}  The precision of the control's calls rises with the horizon, from
$0.27$ at one minute, where more than half of the calls end with no move, to $0.56$ at
twelve hours, and the share of calls that end on the wrong side rises with it, from $0.15$
to $0.33$; the moves when right and when wrong have about the same size, twenty to forty
ticks on average, far above the one-tick median move, because the mean is made by the
wide-spread markets that move in jumps.  The gross edge is two ticks per share at one minute,
four and a half at fifteen minutes and eight to nine from one hour on.  Against it, the
median spread at the calls is twenty to thirty ticks and the spread actually paid by a
taker, dominated by the same wide markets, is eighty to one hundred; the fee is under one
tick for the two legs on average, six to eight on the crypto markets that charge it.  The
necessary condition of \cref{eq:econ:taker} fails on the typical market at every horizon, for
every model and for the oracle too: even a model that is never wrong has a gross edge of
eighteen to forty ticks, the oracle's $\mu^{+}(1 - z)$, which is a fifth of the spread it
would pay.

\paragraph{The taker.}  \Cref{tab:econ:pnl} says the same in money.  Executed as a taker, the
calls of the control lose sixty-eight to one hundred ticks per share at every horizon; removing
the fee changes the loss by less than one tick; the oracle loses forty to ninety.  The loss
shrinks with the horizon because the spread is the same while the moves are larger, but it
never closes.  The one place where a taker breaks even is the books with a spread of at most
two ticks, a tenth of the calls: there the edge of two to three ticks pays the spread and the
fee almost exactly, between $-0.9$ and $+0.9$ ticks per share, and only the most confident
decile of those calls is clearly positive, from under one tick at one minute to twenty or
thirty ticks per share at fifteen minutes and one hour, on one percent of the samples.  The spread, not the fee, is what eats the profit of a
direction call: on this venue in March 2026 a taker pays the fee on a minority of markets
and the spread on all of them, and the spread is a hundred times larger.

\paragraph{The maker.}  Resting at the best quote changes the sign.  A maker--taker is
filled on $7\%$ of the calls at one minute and on half of them at twelve hours, since a
resting order has more time to be reached, and the filled calls earn $-4$ ticks per share at
one minute, nothing at fifteen minutes and ten to fourteen ticks from one hour to twelve.
The one-minute loss is adverse selection in its plain form: the calls that get filled are
those whose price came to the order, and when they are wrong they are wrong by thirty-seven
ticks against nineteen on all calls; from one hour on the selection works the other way,
because the fills concentrate on the active markets where the control is more often right
($0.52$ to $0.57$ precision on the filled calls against $0.46$ to $0.56$ on all).  On the
most confident decile of its calls the control earns forty-six to seventy-nine ticks per
filled share from one hour on, with a precision of $0.76$ to $0.77$, and the oracle earns
forty to eighty: the model's maker gain on all calls is a tenth to a third of what perfect
foresight would earn, and on its best decile it is close to it.  The maker--maker bound adds
the exit half spread, thirty ticks, and shows what the second fill would be worth: more than
the whole edge of any model.  Which side is better is therefore not in doubt: per filled
share a maker is eighty to one hundred ticks ahead of a taker on all calls at every horizon,
and the costs that remain to it are the fill, which the simulation grants optimistically,
and the adverse selection, which it measures.

\paragraph{Risk, horizons and kinds of market.}  The per-share figures hide a wide
dispersion: the standard deviation of the profit of a filled maker--taker trade is 50 to 145
ticks against a mean of $-4$ to $+14$, so that a single trade is a coin toss with a slight
tilt (hit rates of $0.45$ to $0.61$) and the mean is only visible over hundreds of thousands
of trades; per unit of capital at risk, the aggregate return of the maker--taker on all
calls is $1.5$ to $2\%$ per holding period from one hour on, and negative at one minute.
The economics favour the long horizons, where the same spread is paid against an edge that
is five times larger and a resting order is filled seven times as often as at one minute.  By kind of market, the
maker--taker at one hour is positive in every category, most on esports and sports
($+28$ and $+16$ ticks per filled share, the in-play markets where the moves are large and
the fills frequent) and least on economy and technology and on crypto ($+1$ and $+2$, where
the fee of five ticks takes most of the edge); at fifteen minutes it is negative on sports,
crypto and economy and positive on esports, weather and politics.  What the simulation does
not include bounds its reading: our own orders would move the book and lengthen the queue,
a real fill comes later than the first print at our price, the exit of the maker--maker
bound is not guaranteed, and the fee schedule of the venue has since been extended to most
categories at rates between $1\%$ and $5\%$, which at $p = 0.5$ is $2.5$ to $12.5$ ticks per
leg, still below the typical spread but above the edge of any model on the tight books.
The simulation fixes orders of magnitude, not a trading record: a direction call on this
venue is worth a few ticks per share, it cannot be collected by crossing the spread, and
a resting order collects part of it at the cost of a fill that is not guaranteed.
